% !TeX root = ../../main.tex
\section{Introduction}
\label{sec:mobiwac:intro}

Location-based social networks, such as Gowalla~\cite{cho2011gowalla}, let people share the places that they visit as check-ins, which together record how people move through a city. Individual mobility is highly predictable in principle~\cite{song2010limits}, so a service that anticipates the next move can prepare ahead of time, caching content where a user is heading~\cite{bastug2014edge} or provisioning capacity at the destination before demand arrives. Handover predictions already let cellular services adapt in advance at the network level~\cite{vielhaus2022handover}; at the city scale, check-in mobility analysis is likewise a design input for mobility-aware services~\cite{moura2025mobilityaware}.

Predicting the next place exactly, a single point of interest (POI), is hard and often more than a service needs. Two coarser questions are usually enough: the next category, the type of place such as food or shopping, and the next region, the part of the city that a user is heading to (a neighborhood-scale unit such as the U.S. census tract). The first captures intent, the second captures geography; we study whether one model should learn both at once.

Sharing one representation across tasks has a cost: in multitask learning (MTL), one model does several jobs at once by sharing most of its parts, so the shared parameters can converge to a compromise optimal for neither task, helping one while hurting the other~\cite{caruana1997multitask}. Our earlier
work reported no consistent multitask advantage for the paired category tasks and attributed it, in part, to this effect~\cite{silva2025mtlnet}; the useful question is where sharing helps, what it costs, and how to share so the gains hold and the cost stays small.

We propose two enhancements. First, we build a check-in-level representation: instead of one fixed vector per place, each check-in gets its own vector that holds its context (the time, nearby places, and recent visits).
This adds a fourth level, the check-in, beneath the place, region, and city levels of hierarchical graph infomax~\cite{huang2023hgi,velickovic2019deep} (Figure~\ref{fig:mobiwac:dataflow}). Second, we train one model that
predicts the next category and the next region in a single forward pass, with a shared trunk (a cross-attention stack where the two tasks exchange
semantic context) and a private spatial path for the region task (Figure~\ref{fig:mobiwac:model}). We evaluate on two settings chosen to differ: five U.S. states of different sizes (Gowalla) and one non-U.S. city (Istanbul).\footnote{\label{fn:mobiwac:code}Code (model, representation, baselines, statistical tests): \url{https://github.com/VitorHugoOli/PoiMtlNet/tree/mobiwac}. Both data sources are public: the category-annotated Gowalla dump (\url{https://figshare.com/articles/dataset/gowalla_data/22126586}) and Massive-STEPS~\cite{wongso2025massivesteps}.}

Our contributions are as follows.
\begin{itemize}
  \item A check-in-level representation that extends hierarchical graph infomax from the place to the individual visit, improving next-category prediction over a standard place embedding across all six datasets, by $+0.23$ to $+6.29$ points of macro-F1 (a balanced score that counts every
    category equally) (Table~\ref{tab:mobiwac:representation}). The direction is
    consistent even where the size is small, because a single fixed vector
    per place cannot separate places that serve several purposes, which
    per-visit context recovers.
  \item A single model for joint next-category and next-region prediction, sharing semantic context across tasks while keeping a private spatial path for the region task; to our knowledge, the first work to treat fine-grained region as an end target of equal standing (Section~\ref{sec:mobiwac:related-tasks}). In one forward pass, it outperforms a dedicated
    single-task region model at Texas and California, the two datasets with the
    largest region vocabularies, and at the other four it stays within a
    two-point margin registered before any result was read (statistical
    non-inferiority, assessed with two one-sided tests). On category it
    outperforms the dedicated model at Florida, and the five remaining
    differences are equivalent to zero within half a point, even though each
    dedicated model is tuned per dataset (Table~\ref{tab:mobiwac:results}).
  \item An empirical account, across five U.S. states and one non-U.S. city, of where the joint
    region gain appears. The two datasets where the joint model outperforms are the two with the
    largest region vocabularies, Texas and California, while the four with smaller vocabularies sit
    inside the two-point margin (Figure~\ref{fig:mobiwac:deltas}). The grouping is reported as an
    observation rather than a law: the ordering does not hold inside the pair, since California has
    more regions than Texas and a slightly smaller gain, and region count co-varies with corpus
    size across these states. All joint and dedicated results average
    four random initializations (Section~\ref{sec:mobiwac:results-part2}).
\end{itemize}

The remainder of this chapter presents related work (Section~\ref{sec:mobiwac:related}), the problem (Section~\ref{sec:mobiwac:problem}), the method and setup (Sections~\ref{sec:mobiwac:method} and~\ref{sec:mobiwac:setup}), results (Section~\ref{sec:mobiwac:results}), and discussion and conclusions (Sections~\ref{sec:mobiwac:discussion} and~\ref{sec:mobiwac:conclusion}).
